%!TEX root = ../main.tex
\section{Second-Order Memory over the Read Path}
\label{sec:secondorder}

The framework of Section~\ref{sec:framework} evaluates each query in isolation. The hundredth query for a user is therefore served in the same way as the first. Prior memory systems share this limitation. They observe only the \emph{write path}, which sessions arrived, and never the \emph{read path}, which stored records answered which queries. \sys{} closes this loop with a per-user workload log $\mathcal{W}_u$ (Section~\ref{subsec:log}) and a feedback rule that reshapes the catalog's physical organization (Section~\ref{subsec:feedback}).

\subsection{Recording the Read Path}
\label{subsec:log}

Adapting storage to a user presupposes a record of what that user's queries actually consumed. Prior systems never keep that record. This subsection specifies what \sys{} logs, and it keeps the log small enough that logging does not itself become a scaling problem. For each user $u$, \sys{} maintains a compact log of query outcomes
\[
\mathcal{W}_u \;=\; \{\,(q_t,\ \rho(q_t),\ E_t^{+},\ \tau_t)\,\}_{t\leq T},
\]
where $E_t^{+}\subseteq\mathcal{E}$ is the set of catalog entities that appeared in the accepted answer $\hat{a}_t$ or its supporting evidence. The log carries three signals. The sequence of plan classes $\{\rho(q_t)\}$ describes \emph{what the user asks}. The sets $\{E_t^{+}\}$ record \emph{what memory answers it}. Sessions left unread beyond a staleness horizon indicate \emph{what memory is dead weight}. After Alice's first six queries, her log records a \point-heavy class mix over financial and scheduling topics, hits on $\textsf{Raiffeisen Basel}$ and her calendar events, and no reads on her older travel sessions.

\stitle{Per-Entity Statistics.}
The log is statistical, not verbatim. For each catalog entity $e$ we keep three running numbers only, an issue count $n(e)=|\{t\leq T: e\in q_t\}|$ (how often $e$ has been mentioned by any query), an answer-hit count $h(e)=|\{t\leq T: e\in E_t^{+}\}|$ (how often $e$ appeared in an accepted answer), and a last-hit time $\tau^{\star}(e)=\max\{\tau_t: e\in E_t^{+}\}$. Storage is therefore $O(|\mathcal{E}|)$ per user and independent of query volume $T$. A DBMS keeps histograms and access counters instead of a replay log of every statement, for the same reason, since the optimizer needs the workload's distribution and not its transcript.

\stitle{What Counts as an Accepted Answer.}
No gold labels exist at deployment time, and acceptance must therefore be defined operationally. \sys{} marks an answer accepted when the executor commits to it under Eq.~(\ref{eq:short}), which happens when two independent samples agree, and records only the entities the operator read in producing it. The score moves a ranking and never a decision, and a promoted entity changes which records the operator sees first, never what it computes or may return. The loop is a retrieval prior, not a cache of asserted facts.

\stitle{Per-User Isolation.}
$\mathcal{W}_u$ and every quantity derived from it are strictly per-user, with no cross-user aggregation. A new user starts from a neutral prior and adapts as queries arrive.

\subsection{Reshaping Storage from Observed Workload}
\label{subsec:feedback}

A log is useful only if something reads it. This subsection turns the log into a single per-entity score and then spends that score on three decisions about where memory physically sits, leaving query semantics untouched throughout.

\stitle{Workload-Driven Priority Score.}
Each catalog entity $e$ carries a priority score
\begin{equation}
\label{eq:pi}
\pi(e) \;=\; \underbrace{\alpha\cdot h_r(e)}_{\text{used in answers}}\;+\;\underbrace{\beta\cdot r(e)}_{\text{used recently}}\;+\;\underbrace{\gamma\cdot c(e)}_{\text{matches the class}},
\end{equation}
where $h_r(e)=h(e)/n(e)$ is the hit rate on queries that mention $e$, $r(e)=\exp\bigl(-\lambda(\tau_{\mathrm{now}}-\tau^{\star}(e))\bigr)$ is a recency term with decay rate $\lambda$, and $c(e)$ is the fraction of queries mentioning $e$ whose plan class matches $e$'s record type. The three weights are fixed at $(\alpha,\beta,\gamma)=(0.5,0.3,0.2)$ for every benchmark and backbone we report, with direct workload evidence given the dominant share and the remaining two terms acting as tie-breakers. We do not tune them per dataset, and Section~\ref{subsec:exp5} reports the effect of removing each term outright.

\stitle{Storage-Level Actions.}
$\pi(e)$ is refreshed after every query and drives the physical organization of the catalog through three actions, governed by fixed system-wide thresholds. \emph{Index warming} keeps the highest-priority entities by $\pi$ in a resident cache that biases candidate retrieval. \emph{Grounding-block composition} assembles the evidence block prepended to an operator's prompt from the top candidate entities by $\pi(e)$, and not from a static set. \emph{Session eviction} compresses or drops any session $s_i$ whose entities all fall below a cold threshold $\pi_{\mathrm{cold}}$. By Alice's seventh query, her calendar events sit at the top of the grounding block while the never-read travel sessions have been compressed away. When \opAGG{} then enumerates candidates for the outdoor question the eight relevant events are the first records it reads. The operator itself is unchanged, and the workload has reorganized only the catalog it reads.

All three actions operate on \emph{where} memory sits, not on \emph{what} a query means. The access paths $\mathcal{I}_e,\mathcal{I}_a$ simply return workload-reweighted results, and every physical operator benefits without issuing an additional LLM call.

\stitle{Remarks.}
The priority score of Eq.~(\ref{eq:pi}) implements the self-tuning-index principle~\cite{chaudhuri1997autotuning,chaudhuri2007selftuning} and the workload-adaptive stance of learned optimizers~\cite{neoopt,baoopt,balsa}, transferred to agent memory, where the system observes the query log, updates statistics online, reshapes physical storage, and leaves query semantics untouched. In \sys{} the recommender is $\pi(e)$, the log is $\mathcal{W}_u$, and the storage is the catalog. To our knowledge, no prior agent memory system closes this workload-to-storage loop. This is why prior accuracy on a given user stays flat as that user issues more queries (Section~\ref{subsec:exp6}).
